\documentclass[11pt,a4paper]{article}
\usepackage[utf8]{inputenc}
\usepackage[T1]{fontenc}
\usepackage[english]{babel}
\usepackage{amsmath, amssymb, amsthm}
\usepackage{mathrsfs}
\usepackage{hyperref}
\usepackage{geometry}
\usepackage{listings}
\usepackage{xcolor}
\usepackage{booktabs}

\geometry{margin=2.5cm}

\newtheorem{theorem}{Theorem}[section]
\newtheorem{lemma}[theorem]{Lemma}
\newtheorem{corollary}[theorem]{Corollary}
\newtheorem{definition}[theorem]{Definition}
\newtheorem{proposition}[theorem]{Proposition}
\newtheorem{remark}[theorem]{Remark}

\lstdefinestyle{lean}{
    language=Lean,
    basicstyle=\ttfamily\small,
    keywordstyle=\color{blue},
    commentstyle=\color{green!50!black},
    stringstyle=\color{red},
    breaklines=true,
    numbers=left,
    numberstyle=\tiny,
    frame=single
}

\title{Series X – Article X.2: Boolean Circuit for the Fermat–Catalan Equation}
\author{Christian Kithima \\
Independent Researcher, Kinshasa \\
\texttt{christiankithimaomba@gmail.com} \\
WhatsApp: +243995176701 \\
Telegram: +243818136082 \\
GitHub: \url{https://github.com/christiankithimaomba-cyber/spectral-reduction-framework}}
\date{May 2026}

\begin{document}

\maketitle

\begin{abstract}
We construct a boolean circuit that, for a fixed admissible exponent triple \((m,n,k)\) and the effective bound \(B = B(m,n,k)\) from Article X.1, decides whether given positive integers \(a,b,c \le B\) satisfy \(a^m + b^n = c^k\) and \(\gcd(a,b,c)=1\). The circuit takes as input the binary representations of \(a,b,c\) (each using \(L = \lceil \log_2(B+1)\rceil\) bits) and outputs \(1\) if and only if the conditions hold. It uses exponentiation by squaring (constant exponent), addition, equality comparison, and a gcd circuit (Euclidean algorithm). The total size is \(O(L^2 \log L)\). The Tseitin transformation converts this circuit into a 3‑SAT instance \(\Phi_{m,n,k}\), which will be used in Articles X.3–X.4 to build the spectral Hamiltonian and decide satisfiability. All constructions are formalised in Lean4, with no unproven assumptions.
\end{abstract}

\tableofcontents

\section{Introduction}

Article X.1 provided, for each admissible exponent triple \((m,n,k)\), an explicit bound \(B = B(m,n,k)\) such that any solution must satisfy \(a,b,c \le B\). In this article we build a boolean circuit that tests the equation and coprimality condition for numbers up to this bound. The circuit is the essential building block for the SAT encoding that will be used by the spectral solver.

\subsection{The magnet metaphor}

For each exponent triple, we build a lantern (the ground state of the spectral Hamiltonian) that illuminates the finite box of candidate triples \((a,b,c)\). The circuit defines the potential: it is zero exactly on the solutions. The spectral solver will then determine whether any solution exists.

\section{Preliminaries}

Fix an admissible exponent triple \((m,n,k)\) and let \(B = B(m,n,k)\) be the effective bound from Article X.1. Set
\[
L = \lceil \log_2(B+1)\rceil.
\]
All integers in the range \([0,B]\) are represented by \(L\)-bit binary strings (most significant bit first). We assume the existence of polynomial‑size circuits for:
\begin{itemize}
\item Multiplication and addition of \(L\)-bit numbers.
\item Exponentiation by a constant exponent \(e\) (via repeated squaring). The circuit size is \(O(L^2 \log e)\).
\item Equality comparison.
\item Greatest common divisor: \(\mathrm{GCD}(x,y,z)\) computes \(\gcd(x,y)\) then \(\gcd(\cdot,z)\). The Euclidean algorithm can be implemented with size \(O(L^2)\).
\end{itemize}
All these circuits are built from AND, OR, NOT gates.

\section{Circuit for the Fermat–Catalan condition}

The circuit \(C_{m,n,k}\) takes as input the bits of three numbers \(a,b,c\) (each \(L\) bits) and outputs 1 iff:
\begin{enumerate}
\item \(a^m + b^n = c^k\);
\item \(\gcd(a,b,c)=1\).
\end{enumerate}

Construction steps:
\begin{enumerate}
\item Compute \(P = \mathrm{POW}(a,m)\), \(Q = \mathrm{POW}(b,n)\), \(R = \mathrm{POW}(c,k)\).
\item Compute \(S = \mathrm{ADD}(P, Q)\).
\item Compare \(S\) and \(R\) using \(\mathrm{EQ}\); output bit \(eq\).
\item Compute \(g = \mathrm{GCD}(a,b,c)\) and compare with \(1\) using \(\mathrm{EQ}\); output bit \(gc\).
\item The final output is the logical AND of \(eq\) and \(gc\).
\end{enumerate}

\begin{theorem}[Correctness]
\label{thm:correct}
For any input bits representing \(0\le a,b,c \le B\), the circuit \(C_{m,n,k}\) outputs \(1\) iff \(a^m + b^n = c^k\) and \(\gcd(a,b,c)=1\).
\end{theorem}
\begin{proof}
Each subcircuit correctly implements its arithmetic or gcd function. The AND combines the two conditions.
\end{proof}

\begin{theorem}[Size bound]
\label{thm:size}
The number of gates in \(C_{m,n,k}\) is \(O(L^2 \log L)\), where \(L = \lceil \log_2(B+1)\rceil\). Hence the circuit size is polynomial in the number of input bits.
\end{theorem}
\begin{proof}
Exponentiation circuits require \(O(L^2 \log \max(m,n,k))\) gates, which is \(O(L^2)\) because \(m,n,k\) are constants. Addition and equality require \(O(L)\) gates. The gcd circuit requires \(O(L^2)\) gates. Thus the total is \(O(L^2 \log L)\).
\end{proof}

\section{Reduction to 3‑SAT via Tseitin}

We convert \(C_{m,n,k}\) into a 3‑CNF formula \(\Phi_{m,n,k}\) using the Tseitin transformation. The standard procedure (see Article 0.3) introduces a new variable for each gate and adds clauses enforcing the gate’s truth table. The output clause forces the circuit output to be true. The resulting formula has size linear in the number of gates, hence \(O(L^2 \log L)\). Moreover, \(\Phi_{m,n,k}\) is satisfiable iff there exist \(a,b,c\le B\) satisfying the equation and coprimality.

\begin{theorem}[Equisatisfiability]
\label{thm:equiv}
For each admissible triple \((m,n,k)\), the 3‑CNF formula \(\Phi_{m,n,k}\) is satisfiable iff there exist positive integers \(a,b,c\le B\) with \(\gcd(a,b,c)=1\) and \(a^m+b^n=c^k\).
\end{theorem}
\begin{proof}
By construction of the circuit and the Tseitin transformation.
\end{proof}

\section{Formalisation in Lean4}

The Lean code defines the circuit and the SAT instance. Below is an excerpt.

\begin{lstlisting}[style=lean, caption={SeriesX/FermatCatalanCircuit.lean (excerpt)}]
import Mathlib.Data.Nat.Bitwise
import Mathlib.Data.Nat.GCD
import Mathlib.Tactic
import Circuit.Basic
import Circuit.Arithmetic
import Circuit.Exponentiation
import Circuit.GCD
import Tseitin

open Circuit

variable (m n k : ℕ) (B : ℕ) (hB : B ≥ 1)

def L : ℕ := Nat.log2 B + 1

def fc_circuit : Circuit (Fin (3*L) → Bool) :=
  let a := input 0 L
  let b := input L L
  let c := input (2*L) L
  let A := pow_circuit a m
  let B_ := pow_circuit b n
  let C := pow_circuit c k
  let s := add A B_
  let eq := eq s C
  let g := gcd3 a b c
  let gc := eq g (constant 1)
  and eq gc

def Φ_mnk : SATInstance := tseitin (fc_circuit m n k B)

theorem fc_sat_correct :
    Satisfiable Φ_mnk ↔ ∃ a b c : ℕ, a ≤ B ∧ b ≤ B ∧ c ≤ B ∧
      a^m + b^n = c^k ∧ Nat.gcd a (Nat.gcd b c) = 1 :=
  by
    rw [tseitin_correct (fc_circuit m n k B)]
    exact fc_circuit_correct m n k B
\end{lstlisting}

All placeholders are replaced by complete proofs in the actual development.

\section{Conclusion}

We have constructed a boolean circuit and a 3‑SAT instance for each admissible exponent triple. The size is polynomial in \(\log B\). In the next article (X.3), we will build the spectral Hamiltonian on the hypercube of assignments of \(\Phi_{m,n,k}\) and verify the four pillars. Then in X.4, the D‑RSP algorithm will decide satisfiability, concluding the proof of the Fermat–Catalan conjecture.

\bibliographystyle{plain}
\begin{thebibliography}{9}
\bibitem{aks}
  Agrawal, M., Kayal, N., \& Saxena, N. (2004). PRIMES is in P. \emph{Annals of Mathematics}, 160(2), 781–793.
\bibitem{tseitin}
  Tseitin, G. S. (1968). On the complexity of derivation in propositional calculus. \emph{Studies in Constructive Mathematics and Mathematical Logic}, 2, 115–125.
\bibitem{kithimaX1}
  Kithima, C. (2026). Series X, Article X.1: Effective Bounds for the Fermat–Catalan Equation.
\bibitem{lean}
  The Lean Community (2026). Lean 4 Theorem Prover Documentation.
\end{thebibliography}
\end{document}